%Chargement des paquets
\usepackage{amsmath}
\usepackage{amsfonts}
\usepackage{amssymb}
\usepackage{enumerate}
\usepackage{mathtools}
\usepackage{bbm}
\usepackage{xparse, etoolbox}
\usepackage{enumerate}
\usepackage{enumitem}
\usepackage{mathabx}
\usepackage{babel}[french]

%environment
\usepackage{amsthm}
\usepackage{keytheorems}
\usepackage{hyperref} 

%Interface théorème
\newkeytheorem{lem}[
	name = Lemme
]

\newkeytheorem{prop}[
	name = Proposition
]

\newkeytheorem{thm}[
	name = Théorème
]

\newkeytheorem{coro}[
	name = Corollaire
]

\renewcommand*{\proofname}{Démonstration}

\theoremstyle{definition}
\newtheorem{defn}{Définition}

\theoremstyle{definition}
\newtheorem*{nota}{Notation}

\theoremstyle{definition}
\newtheorem*{limi}{Liminaire}

\theoremstyle{remark}
\newtheorem{rmq}{Remarque}

\theoremstyle{plain}
\newtheorem*{fait}{Fait}


%Ensembles de nombres
\newcommand{\N} {\mathbb{N}}
\newcommand{\Ne}{\N^\ast}
\newcommand{\Z} {\mathbb{Z}}
\newcommand{\Q} {\mathbb{Q}}
\newcommand{\R} {\mathbb{R}}
\newcommand{\Rb}{\overline{\mathbb{R}}}
\newcommand{\Rp}{\R_+}
\newcommand{\Rm}{\R_-}
\newcommand{\K} {\mathbb{K}}
\newcommand{\Ket} {\mathbb{K^{\ast}}}
\newcommand{\Cx}{\mathbb{C}}

%Ensembles, tribus
\newcommand{\A}{\mathbb{A}}
\renewcommand{\P}{\mathcal{P}}
\newcommand{\T}{\mathcal{T}}
\newcommand{\F}{\mathcal{F}}
\newcommand{\C} {\mathcal{C}}
\newcommand{\ouv}{\mathcal{O}}
\newcommand{\B}{\mathcal{B}}
\newcommand{\M}{\mathcal{M}}
\newcommand{\etag}{\mathcal{E}}
\newcommand{\etagOp}{\etag^{+}(\Omega)}
\newcommand{\etagO}{\etag(\Omega)}
%%Lp avant quotientage
\newcommand{\Lic}{\mathcal{L}^1}
\newcommand{\Lpc}{\mathcal{L}^p}
\newcommand{\Linfc}{\mathcal{L}^{\infty}}
\newcommand{\Lifull}{\Lic(\Omega, \B, m)}
\newcommand{\Lpfull}{\Lpc(\Omega, \B, m)}
\newcommand{\Linffull}{\Linfc(\Omega, \B, m)}
%%Lp après quotientage
\newcommand{\Li}{L^1}
\newcommand{\Ld}{L^2}
\newcommand{\Lp}{L^p}
\newcommand{\Linf}{L^{\infty}}
\newcommand{\Listd}{\Li(\Omega, m)} 
\newcommand{\Ldstd}{\Ld(\Omega, m)} 
\newcommand{\Lpstd}{\Lp(\Omega, m)} 

%Quelques opérateurs
\DeclareMathOperator{\codim}{codim}
\DeclareMathOperator{\rg}{rg}
\DeclareMathOperator{\tr}{tr}
\DeclareMathOperator{\vect}{Vect}
\DeclareMathOperator{\ima}{Im}
\DeclareMathOperator{\id}{Id}
\DeclareMathOperator{\sign}{sign}
\DeclareMathOperator{\ext}{ext}
\newcommand{\ind}{\mathbbm{1}}
\newcommand*{\spr}[2]{\left(\,#1\,\middle|\,#2\,\right)}
\newcommand{\interior}[1]{%
  {\kern0pt#1}^{\mathrm{o}}%
}
\DeclareMathOperator{\card}{card}
\DeclareMathOperator{\ppcm}{ppcm}

%Commandes de pure flemme
\newcommand{\veps}{\varepsilon}
\newcommand{\intab}{\int_{a}^{b}}
\newcommand{\intR}{\int_{-\infty}^{\infty}}
\newcommand{\intinf}{\int_{- \infty}^{+ \infty}}
\newcommand{\equi}{\Leftrightarrow}
\newcommand{\Kev}{$\K$-espace vectoriel }
\newcommand{\Kevs}{$\K$-espaces vectoriels }
\newcommand{\Rev}{$\R$-espace vectoriel }
\newcommand{\Revs}{$\R$-espaces vectoriels }
\newcommand{\Cev}{$\Cx$-espace vectoriel }
\newcommand{\Cevs}{$\Cx$-espaces vectoriels }
\newcommand{\evn}{(E, \norm{.})}
\newcommand{\interf}[2]{[#1,#2]}
\newcommand{\limb}{\lim\limits_}
\newcommand{\lebn}{\lambda_n}
\newcommand{\pp}{\text{~presque partout}}
\newcommand{\derivp}[2]{\frac{\partial #1}{\partial #2}}
\newcommand{\famiu}{(u_i)_{i \in I}}
\newcommand{\sev}{\text{~s.e.v.}}
\newcommand{\Mnp}{M_{n,p}(\K)}
\newcommand{\Mn}{M_{n}(\K)}
\newcommand{\tq}{\text{~t.q.~}}
\newcommand{\mq}{~montrer~que~}
\newcommand{\et}{\quad \text{et} \quad}
\newcommand{\ou}{\text{~ou~}}
\newcommand{\Kx}{\K[X]}


%Commandes de gars chelou
\renewcommand{\subset}{\subseteq}

%Commande restriction, ty duckduckgo
\def\restr#1#2{\mathchoice
              {\setbox1\hbox{${\displaystyle #1}_{\scriptstyle #2}$}
              \restraux{#1}{#2}}
              {\setbox1\hbox{${\textstyle #1}_{\scriptstyle #2}$}
              \restraux{#1}{#2}}
              {\setbox1\hbox{${\scriptstyle #1}_{\scriptscriptstyle #2}$}
              \restraux{#1}{#2}}
              {\setbox1\hbox{${\scriptscriptstyle #1}_{\scriptscriptstyle #2}$}
              \restraux{#1}{#2}}}
\def\restraux#1#2{{#1\,\smash{\vrule height .8\ht1 depth .85\dp1}}_{\,#2}} 

%Quelques normes
\DeclarePairedDelimiter\abs{\lvert}{\rvert}
\DeclarePairedDelimiter\labs{\bigl\lvert}{\bigr\rvert}
\DeclarePairedDelimiter\norm{\lVert}{\rVert}
\DeclarePairedDelimiterXPP\normi[1]{}\lVert\rVert{_1}{\ifblank{#1}{\:\cdot\:}{#1}}
\DeclarePairedDelimiterXPP\normd[1]{}\lVert\rVert{_2}{\ifblank{#1}{\:\cdot\:}{#1}}
\DeclarePairedDelimiterXPP\normp[1]{}\lVert\rVert{_p}{\ifblank{#1}{\:\cdot\:}{#1}}
\DeclarePairedDelimiterXPP\normq[1]{}\lVert\rVert{_q}{\ifblank{#1}{\:\cdot\:}{#1}}
\DeclarePairedDelimiterXPP\norminf[1]{}\lVert\rVert{_\infty}{\ifblank{#1}{\:\cdot\:}{#1}}

%Bracket en angle
\DeclarePairedDelimiter\angl{\langle}{\rangle}

%Set, parenthèses
\newcommand{\set}[1]{\{ #1 \}}
\newcommand{\bigset}[1]{\bigl\{ #1 \bigr\}}
\newcommand{\Bigset}[1]{\Bigl\{ #1 \Bigr\}}
\newcommand{\bigpar}[1]{\bigl( #1 \bigr)}
\newcommand{\Bigpar}[1]{\Bigl( #1 \Bigr)}

%Opitmisation des opérateurs de comparaison
\DeclarePairedDelimiter\tio{o (}{)}
\DeclarePairedDelimiter\gro{O (}{)}


\pagestyle{fancy}

\fancyhead[C]{}
\fancyhead[R]{}

\begin{document}

\underline{Source :} Gourdon - Analyse - pages 219 et 130

\underline{Leçons :}
\begin{itemize}
\item 223 : Suites num´eriques. Convergence, valeurs d’adh´erence. Exemples et applications.
\item 224 : Exemples de d´eveloppements asymptotiques de suites et de fonctions.
\item 236 : Illustrer par des exemples quelques m´ethodes de calcul d’int´egrales de fonctions d’une ou plusieurs variables.
\end{itemize}

\begin{defn}[Intégrale de Wallis]
Pour $n \in \N$ on définit l'intégrale de Wallis de rang $n$ par :
\[ W_n = \int_{0}^{\pi/2} \sin^n{x} dx . \]
\end{defn}

\begin{lem}
\label{lem:1}
Pour tout $p \in \N$, on a :
\[ W_{2p+1} = \dfrac{4^p (p!)^2}{(2p+1)!} \et W_{2p} = \dfrac{2p!} {4^p (p!)^2} \dfrac{\pi}{2} . \]
\end{lem}

\begin{proof}
On remarque que $W_0 = \pi/2$ et $W_1 = 1$. \\
Pour $n \geq 2$, on réalise une intégration par partie en posant :
\[ \begin{matrix}
u(x) = \sin^{n-1}{x} & u'(x) = (n-1) \sin^{n-2}{x} \cos{x} \\
v'(x) = \sin{x} & v(x) = - \cos{x}
\end{matrix} \]
on a alors :
\begin{align*}
W_n & = \underbrace{ \left[ -\cos{x} \sin^{n-1}{x} \right ]}_{=0} + (n-1) \int_{0}^{\pi/2} \sin^{n-2}{x} \cos^2{x} dx \\
& = (n-1) \int_{0}^{\pi/2} \sin^{n-2}{x} (1-\sin^2{x}) dx \\
& = (n-1) W_{n-2} + (n-1) W_n ,
\end{align*}
d'où on tire la relation de récurrence $W_n = \dfrac{n-1}{n} W_{n-2}$. En réalisant une récurrence sur les intégrales de rangs pairs
et une seconde sur les intégrales de rangs impairs, on trouve :
\[ \forall p \in \N, \quad W_{2p} = \dfrac{(2p-1) \cdots 3 \cdot 1}{2p \cdots 4 \cdot 2} \dfrac{\pi}{2} \et
W_{2p+1} = \dfrac{2p \cdots 4 \cdot 2}{(2p+1) \cdots 3 \cdot 1} . \]
\end{proof}

\begin{lem}[Formule de Wallis]
\label{lem:Wallis}
Lorsque $p \to \infty$, on a :
\[ \dfrac1{p} \left ( \dfrac{2p \cdots 4 \cdot 2}{(2p-1) \cdots 3 \cdot 1} \right )^2 \to \pi \]
\end{lem}

\begin{proof}
Soit $x \in [0, \pi/2]$ et $p \in \Ne$, on a $\sin{x} \in [0,1]$ donc :
\[ \sin^{2p+1}{x} \leq \sin^{2p}{x} \leq \sin^{2p-1}{x} , \]
d'où on tire, par intégration, 
\[ W_{2p+1} \leq W_{2p} \leq W_{2p-1} \qquad \text{donc} \qquad 1 \leq \dfrac{W_{2p}}{W_{2p+1}} \leq \dfrac{W_{2p-1}}{W_{2p+1}}
= \dfrac{2p+1}{2p} , \]
en utilisant le résultat du lemme $\ref{lem:1}$ (avec $W_{2p-1} = W_{2(p-1)+1}$). 
On déduit, par passage à la limite, que :
\[ \lim_{p \to \infty} \dfrac{W_{2p}}{W_{2p+1}} = 1 , \qquad (*) \]
avec :
\[ \dfrac{W_{2p}}{W_{2p+1}} = \dfrac{\pi}2 \left ( \dfrac{(2p-1) \cdots 3 \cdot 1}{2p \cdots 4 \cdot 2} \right )^2 (2p+1) . \]
Autrement dit :
\[ \lim_{p \to \infty} p \left ( \dfrac{(2p-1) \cdots 3 \cdot 1}{2p \cdots 4 \cdot 2} \right )^2 = \dfrac1{\pi} . \]
%d'où on tire l'équivalence pour $p \to \infty$ :
%\[ \dfrac{(2p-1) \cdots 3 \cdot 1}{2p \cdots 4 \cdot 2} \sim \dfrac1{\sqrt{p \pi}} , \]
%soit, toujours d'après le lemme $\ref{lem:1}$ et la relation (*) :
%\[ W_{2p} \sim \dfrac12 \sqrt{\dfrac{\pi}{p}} \et W_{2p} \sim W_{2p+1} . \]
%On conclut en remplaçant $p$ par $n/2$ dans l'équivalence.
\end{proof}

\begin{thm}[Formule de Stirling]
Lorsque $n \to \infty$, on a :
\[ n! \sim \sqrt{2\pi n}\cdot n^ne^{-n} . \]
\end{thm}

\begin{proof}
On définit, pour $n \in \Ne$, les suites $(u_n)$ et $(v_n)$ par :
\[ u_n = \dfrac{n^ne^{-n}\sqrt{n}}{n!} \et v_n=\log{u_{n+1}/u_n} . \]
On estime $v_n$ pour $n \to \infty$
\begin{align*}
v_n & = \log \left [ \left ( \dfrac{n+1}{n} \right )^{n+1/2} e^{-1} \right ] \\
& = -1 + \left ( n+\dfrac12 \right ) \left ( \dfrac1{n} - \dfrac1{2n^2} + O \left ( \dfrac1{n^3} \right ) \right ) \\
& = -1 + 1 -  \dfrac1{2n} + \dfrac1{2n} + O  \left ( \dfrac1{n^2} \right ) \\
& = O  \left ( \dfrac1{n^2} \right )
\end{align*}
Ainsi, la suite $(\abs{v_n})$ est dominée par la suite $(1/n^2)$ et la série $\sum 1/n^2$ est convergente, on en déduit que $\sum v_n$
converge. Or, pour $N \in \Ne$ :
\begin{align*}
\sum_{n=1}^{N} v_n & =  \sum_{n=1}^{N} ( \log{u_{n+1}} - \log{u_n} ) \\
& = \log{u_{N+1}} - \log{u_1}
\end{align*}
Par passage la limite et par continuité de la fonction $\log$ on en déduit que la suite $(u_n)$ admet une limite en l'infini.
On note $k = \lim{u_n} > 0$ ($k=0$ est impossible car $\log{u_n}$ converge vers une limite finie) et on déduit que :
\begin{align} n! \sim k \sqrt{n} n^n e^{-n} \qquad n \to \infty . \label{eq:1} \end{align}
D'après la formule de Wallis, on a :
\[ \lim_{p \to \infty} \dfrac1{p} \left ( \dfrac{2p \cdots 4 \cdot 2}{(2p-1) \cdots 3 \cdot 1} \right )^2 = \pi \]
ce qu'on peut réécrire ainsi :
\[ \lim_{p \to \infty} \dfrac1{p} \left ( \dfrac{2^{2p} (p!)^2}{(2p)!} \right )^2 = \pi , \]
soit, en utilisant l'équivalent (\ref{eq:1}) : 
\[ \dfrac1{p} \left ( \dfrac{2^{2p} (p!)^2}{(2p)!} \right )^2 \sim \dfrac{2^{4p}k^4p^{4p+2}e^{-4p}}{pk^2(2p)^{4p+1}e^{-4p}}
= \dfrac{k^2}2 , \]
autrement dit $\pi = k^2/2$ soit $k=\sqrt{2\pi}$, ce qui permet de conclure.
\end{proof}

\end{document}